\section{Intermediate terminal orders}\label{app:middle-tail}

\begin{proposition}\label{prop:middle-tail}
Assume the conditions of Lemma~\ref{lem:support-reduction}.
For \(2200\le N<e^{15}\), a common factor requires
\(\alpha<.03\) and \(b>N\). For \(3150\le N<e^{15}\)
there is no common factor.
\end{proposition}

Write \(t=\log(1/\alpha)\), \(V=\log(M+1)\) and
\(\sigma=2U+2\), where \(U=\log N\).
From Lemma~\ref{lem:gap-caps} we may use \(d\le233\) below
\(20{,}000\) and \(d\le565\) throughout the appendix.
If \(X\) bounds the prime divisors of the lifted norm \(K\), the
right side of \eqref{eq:pade-norm} is at most
\begin{equation}\label{eq:intermediate-prime-head}
 S(X)+\pi(X)V/N,\qquad S(X)=\sum_{p\le X}\frac{\log p}{p-1}.
\end{equation}
In the applications below \(V/\log X<j+.001\) whenever
\(M\le X^j\), since \(M\ge7N\) and \(X\ge9\).
The baseline from Lemma~\ref{lem:resultant} is
\(\log\Lambda> U-1.011-g\). The weaker constant \(1.016\)
also suffices for the estimates of proper cosets.

\begin{lemma}\label{lem:small-prime-sums}
For each of the pairs
\[
 (x_0,c)=(21,.22),(34,.30),(43,.33),(100,.39),
          (117,.425),(137,.449),(150,.45),(343,.49)
\]
one has \(S(x)<\log x-c\) for \(x\ge x_0\).
Also \(\pi(x)\log x/x<1.224\) for \(x\ge117\).
\end{lemma}

\begin{proof}
For \(S(x)-\log x\) below \(3000\), check the listed left endpoints
and the primes, since the function decreases between primes.
For the prime count, check \(117\) and the primes up to \(300\),
since \(\log x/x\) decreases for \(x>e\).
After these ranges the bounds in Lemma~\ref{lem:prime-estimates}
suffice. The supplement uses directed interval arithmetic for the
finite comparisons.
\end{proof}

\subsection{Proper rotation cosets}

Let \(k_0\) be the size of each rotation coset in
Lemma~\ref{lem:kummer}. Suppose \(q=d/k_0\ge2\), and put
\(P=M^{1/(q+1)}\). Then all norm primes are at most \(P\)
and \(4d\le2H<P\). In particular,
\begin{equation}\label{eq:proper-coset-two-caps}
 P^q<\min\{N/(4\alpha),N^2/4\}.
\end{equation}
The first inequality uses \(M=dN/\alpha<PN/(4\alpha)\),
and the second uses \(M<dN^2\).

If \(k_0=1\) and \(d\ge3\), the inequality
\((4d)^{d+1}<dN^2\) gives \(d<U\). Indeed
\(d\log(4d)-\log d\) increases for \(d\ge3\) and exceeds
\(2U\) at \(d=U\) for \(U\ge\log2200\).
Use \(X=(UN^2)^{1/4}>P\) and \(s<3U/4\) in the
baseline with \(g\le7s/32\). The resulting surplus over
\eqref{eq:intermediate-prime-head} is at least
\[
 U-1.016-\frac{21}{128}U
 -\log X+.22-4.001\cdot1.256\,X/N.
\]
This exceeds \(1.10\) at \(N=2200\) and increases thereafter.
If \(k_0=1,d=2\), use \(X=(2N^2)^{1/3}\), \(s<1.5\) and
the prime bound with \(c=.33\) in place of \(.22\).
The surplus exceeds \(.95\) at \(N=2200\) and increases.

Suppose \(k_0\ge2\) and \(\alpha\ge.03\). For
\(\alpha\in[\ell,u]\) put \(X=\sqrt{N/(4\ell)}\).
By \eqref{eq:proper-coset-two-caps}, \(X>P\) and \(M<X^3\).
At least a third of each coset lies in the left half-plane.
The one-sided gain from Lemma~\ref{lem:resultant} gives the
lower bound
\[
 U-.006-\tfrac23(1.01+\kappa\sigma)
       +\tfrac13\log((1-u)/u)
\]
on \(\log\Lambda\), using \(g\le\kappa s\).
The upper bound from \eqref{eq:intermediate-prime-head} is
\(\log X+c+3.001\cdot1.256\,X/N\).
The following lower margins at \(N=2200\) suffice since the
gain grows faster than \(\log X\) and \(X/N\) decreases.
\[
\begin{array}{cccc}
\toprule
{[\ell,u]}&\kappa&c&\text{surplus exceeds}\\
\midrule
{[.03,.1]}&.134&-.30&1.35\\
{[.1,1/3]}&.160&.20&.75\\
{[1/3,.75]}&.218&.20&.14\\
\bottomrule
\end{array}
\]

For \(\alpha<.03\) and \(q\ge3\), take
\(X=N^{2/3}/4^{1/3}>P\). We have \(M<X^4\).
If \(D\) bounds \(d\), the baseline exceeds
\(U-1.016-.00384D\).
Subtracting the upper bound
\(\log X-.22+4.001(X/N)(1+1.2762/\log X)\)
gives a surplus which increases with \(N\) for a fixed \(D\).
At \(N=2200,D=233\) and \(N=20{,}000,D=565\) it
exceeds \(1.08\) and \(.68\), respectively.

For \(\alpha<.03\) and \(q=2\), use \(X=N/2\), \(M<X^3\)
and \eqref{eq:intermediate-prime-head} to bound the norm
charge by \(U+.862\).
When \(d=4\), half of the size-two coset has gain at least
\(m(1-\alpha)/\alpha\), and the mean log gain exceeds
\(U+1.21\). For \(6\le d\le18\), the small-coset form of
Lemma~\ref{lem:angles}, with \(k_0=d/2\), gives a gain greater
than \(U+1.89\). For \(d\ge20\), use \(g_0\le3.18\),
\(g_0/k_0\le.10868\), \(r/k_0<.111\) and \(J<.153\)
in \eqref{eq:pade-angles}, together with
\(\log C'>t-.7237\) and \(1-1/k_0\ge.9\).
Dropping the positive \(2\log(2k_0/\pi)/k_0\) leaves a gain
exceeding \(U+1.80\).
This proves that the rotation coset is full.

\subsection{Full cosets at moderate and large ratios}

First take \(2200\le N<20{,}000\).
Here \(a\le(97/3)d\), \(d\le233\) and \(H\le467\).
Thus \(aN+H<189{,}695{,}659=X_{222}\).
Lemma~\ref{lem:local-gap} gives \(H<222\) and \(d\le110\).
For \(N\le3500\) these caps give \(aN+H<X_{154}=17{,}051{,}707\),
so \(d\le76\).

On a parameter interval \(\ell\le\alpha\le u\), order
the \(d\) rotations by increasing absolute argument and use
\(c_j=\cos(\pi j/d)\), an upper bound for the cosine of
the argument in position \(j\). To state the pointwise gain,
let \(W=(1-u-c)/u\) for \(c\ge0\).
Substitute \(v=1-e^{-x/d}\) in the exponent of the second
integrand in \eqref{eq:pade-integrals}. Its logarithm divided
by \(m\) is at most \(-Wx+\log(1-e^{-x})\).
For \(W>0\) the maximum is \(-J(W)\), where
\(J(W)=(1+W)\log(1+W)-W\log W\).
Thus a lower bound for the log gain, minus \(\log m\), is
\(\Phi_u(c)=\max\{-1.01+J(W),\log W\}\),
by the same Stirling estimate as in Lemma~\ref{lem:resultant}.
Here \(J(W)<5\), which is absorbed by the Stirling slack at
\(m\ge2199\).
For \(c<0\) use \(W=(1-u)/u\) in the integrand and
replace \(\log W\) in this maximum by \(\log(W(1-c))\).
For \(c\ge0,W\le0\) use the baseline
\(-1.01-\kappa(u)du\), where \(\kappa(u)=g/s\) at \(\alpha=u\)
in \eqref{eq:pade-loss}.
In the smaller terminal band we also use the lower bound
\[
 -1.01-\frac d u
       \left(c-(1-u)-(1-u)\log\frac{c}{1-u}\right)
\]
when this improves the baseline. It follows by discarding
the beta tail factor and maximizing the remaining exponent.
The loss at \(\alpha\) is zero for \(c\le1-\alpha\) and
otherwise the same formula with \(\alpha\) in place of \(u\).
It increases with \(\alpha\), since the derivative has the sign
of \(\log(c/(1-\alpha))-c+(1-\alpha)\ge0\).

By Lemma~\ref{lem:degree-cap}, a prime bound is
\(\max(a+2d-1,\sqrt M)\). It lies between \(X_1\) and \(X_0\),
where
\[
 X_0=\max\{d(1+1/\ell),\sqrt{dN/\ell}\},\qquad
 X_1=\sqrt{\max(7,d/u)N}.
\]
The upper bound for the right side of the norm inequality is therefore
\begin{equation}\label{eq:moderate-norm}
 \log X_0-c_0+2.001(1+1.2762/\log X_1)X_0/N.
\end{equation}
Use \(c_0=.39\) in the smaller band and \(c_0=.45\) in the
larger band. The first term grows with slope at most \(1/2\)
in \(U\) and the last decreases, while the lower bound
\(\log(N-1)+d^{-1}\sum_j\Phi_u(c_j)\) grows with slope greater
than one. Since \(d\ell\le s<2U+2\), compare at
\(N=\max\{N_{\min},\exp((d\ell-2)/2)\}\).
Use the following subdivision of \([.03,.75]\).
Intervals with \(d(1/\ell-1)<5\) contain no eligible support.
The margins are obtained by outward rounding as described in
the supplement.
\[
\begin{array}{rrlc}
\toprule
N_{\min}&d_{\max}&\text{interval widths}&\text{surplus exceeds}\\
\midrule
2200 &76 & .005\text{ to } .05,\ \text{then } .01 &.070\\
3500 &110& .02 &.089\\
\bottomrule
\end{array}
\]

Now take \(20{,}000\le N<e^{15}\). The gap table gives the
successive caps \(D=257,300,565\) at starting orders
\(20{,}000,32{,}768,65{,}536\). In each range \(D/N<.015\),
so \((1+\alpha)^2 b\le D(1+\alpha)^2/\alpha<N\).
Hence \(\sqrt M\) bounds the norm primes. Using
\(S(\sqrt M)<\log\sqrt M-.49\), the right side of the norm inequality is below
\begin{equation}\label{eq:moderate-high-norm}
 \tfrac12(U+\log b)-.49+2.46\sqrt{b/N}.
\end{equation}
At angles with \(\cos\theta\le1-\alpha\) the beta integral
has gain at least \(me^{-1.01}\). The exceptional fraction
is at most \(\arccos(1-\alpha)/\pi+1/d\). For
\(\ell\le\alpha\le u\), use \(\eta\ge\arccos(1-u)/\pi\).
At \(d\ge10\), at least a fraction \(.45\) lies in the left
half-plane, where one may use the greater gain \(m(1-\alpha)/\alpha\).
Thus, with \(S_0=\min(2U+2,Du)\) and
\(b_0=\min((2U+2)/\ell^2,D/\ell)\), the surplus exceeds
\[
 U-1.011-\kappa(\eta S_0+u)
  +.45\max\{0,1.01+\log((1-u)/u)\}
  -\bigl(\tfrac12(U+\log b_0)-.49+2.46\sqrt{b_0/N}\bigr).
\]
The gain grows at rate at least \(1-2\kappa\eta\), the
logarithmic term on the norm side at rate at most
\(1/2+1/(2U+2)\), and the square-root term decreases.
The following minima are taken over the three left endpoints.
For \(d\le9\) the same argument before adding the left half-plane
gain uses \(\kappa=.218,\eta=.42,\ s\le6.75,\ b\le300\),
and gives a positive surplus at \(N=20{,}000\) increasing thereafter.
\[
\begin{array}{cccc}
\toprule
{[\ell,u]}&\kappa&\eta&\text{surplus exceeds}\\
\midrule
{[.03,.036]}&.128&.086&.12\\
{[.036,.045]}&.129&.096&.20\\
{[.045,.07]}&.131&.120&.11\\
{[.07,.11]}&.134&.151&.21\\
{[.11,.17]}&.140&.189&.51\\
{[.17,.28]}&.150&.245&.67\\
{[.28,.45]}&.170&.315&.62\\
{[.45,.60]}&.190&.370&.53\\
{[.60,.75]}&.218&.420&.08\\
\bottomrule
\end{array}
\]


\subsection{Small ratios with the upper exponent at most the order}

Take \(X=1.03\sqrt M\), which is a prime bound by
Lemma~\ref{lem:degree-cap}. Indeed
\(a+2d-1<(1+\alpha)b\le1.03\sqrt M\).
For \(d\ge20\), the angular estimate
\eqref{eq:pade-angles} gives
\[
 \log\Lambda-U>
 .95\bigl(\log(100/3)-.7237\bigr)-.002-.05434
 - .111\bigl(1.01+.00384D+\log2\bigr)-.153
\]
with \(d\le D\). In \eqref{eq:intermediate-prime-head},
use \(X/N\le1.03\) and \(X>1.03\sqrt{20N/.03}\).
At \(N=2200,D=233\) the surplus exceeds \(.17\), and at
\(N=20{,}000,D=565\) it exceeds \(.08\).
The right side of the norm inequality minus \(U\) decreases thereafter for each fixed \(D\).
For \(10\le d\le19\), inserting each \(d\) in
\eqref{eq:pade-angles} gives \(\log\Lambda>U+2.41\),
against a right side of the norm inequality below \(U+2.052\) by \(X>343\).
For \(4\le d\le9\) use the small-coset angular bound. The gain
exceeds \(U+2.24\), against a right side of the norm inequality below \(U+2.177\)
by \(X>137\).

For an antipodal coset with \(d=2\), write \(c=|\cos\theta|\),
\(A_0=(1-\alpha)/\alpha\) and \(B_0=e^{-g_0}\).
The two gains have product at least
\(m^2A_0(1+c)\max\{B_0,A_0-c/\alpha\}\).
On \(0\le c\le1-\alpha(1+B_0)\) this is a concave quadratic
in \(c\) and increases on the rest of its domain. The smaller endpoint gives
the lower bound \(m^2A_0B_0(2-\alpha(1+B_0))\).
In particular \(\log\Lambda>U+t/2-.2\).
For \(d=3\) the small-coset bound gives
\(U+2t/3-.435>U+t/2-.2\).
Put \(\mu=\tfrac12\log(b/N)\).
Since \(X>117\), the right side of the norm inequality is at most
\[
 U+\mu+\log(1.03)-.425+
       2.001\cdot1.03\cdot1.224\,e^\mu.
\]
As \(t=U-\log d+2\mu\), the surplus is minimized at \(\mu=0\)
and exceeds \(.96\) at \(N=2200\).
This proves the first conclusion of Proposition~\ref{prop:middle-tail}.

\subsection{Small ratios with the upper exponent exceeding the order}

It remains to exclude this case for \(N\ge3150\).
The bounds \(H<G,\ d\le D=G/2-1\) are shown in the
following table. In addition to the main gap table, apply
Lemma~\ref{lem:local-gap} successively using \(G=320,292,288,282\)
up to \(N=5000,4150,3467,3295\), respectively.
At every step the preceding caps give
\(aN+H<(5/6)D_{\mathrm{old}}N_{\max}^2+H_{\mathrm{old}}<X_G\).
\[
\begin{array}{rrrccc}
\toprule
N\text{ from}&N\text{ through}&G&A_v&A_c&\text{surplus exceeds}\\
\midrule
3150  &3295   &282 &2.504&2.885&.023\\
3296  &3467   &288 &2.506&2.886&.146\\
3468  &4150   &292 &2.514&2.889&.31\\
4151  &5000   &320 &2.522&2.892&.44\\
5001  &6000   &354 &2.531&2.896&.86\\
6001  &10000  &456 &2.58 &3.01 &.59\\
10001 &19999  &468 &2.58 &2.92 &2.37\\
20000 &32767  &516 &3&3&3.88\\
32768 &65535  &602 &3&3&4.66\\
65536 &3269017&1132&3&3&4.87\\
\bottomrule
\end{array}
\]

Let
\[
 \begin{aligned}
 w&=\log(b/N), & v&=\log M=2U+w, & P&=4M^{1/3},\\
 \zeta&=\log P, & t_D&=U+w-\log D, & \xi&=t_D+.337.
 \end{aligned}
\]
We have \(0<w<\log(5D/6)\) and \(P>\max(830,5G/2)\)
in every row. At \(p>P\) we have \(p>2H+d\) and \(p^3>M\).
Lemma~\ref{lem:two-heights} separates the prime divisors of the
norm there according to whether the two Newton edge heights agree.

The primes with equal heights divide
\(\binom{a+2d-1}{2d-1}\), since they divide a number in
\([a+1,a+2d-1]\) and exceed \(2d-1\).
Using \(\binom{x}{j}\le(ex/j)^j\), their total logarithmic weight
is at most
\[
 2d\bigl(1+\log(a/d+2)-\log2\bigr)<2d(t+.337)\le2D\xi.
\]
For the last step the function \(x(\log(b/x)+.337)\) increases
for \(2\le x\le D\), since \(b>N\) and \(D<N/4\).
Their right side of the norm inequality is at most
\(2D\xi(1/(P-1)+V/(N\zeta))\).

For unequal heights, one of the two windows
\([bm-2H,bm]\), \(m=N-1,N\), contains \(qp^2\)
with \(1\le q\le Q=M/P^2=P/64\).
At fixed \(q\) each window contains at most one such \(p>P\)
by the spacing of the squares. Their digit charge is at most \(2QV/N\).
Put \(Z=b(N-1)-2H\). Since \(Z>.9995M\) and \(Q>12\),
the first sum in \eqref{eq:pade-norm} on these primes is at most
\[
 v\,\frac{\sum_{q\le Q}\sqrt q}{\sqrt Z-\sqrt Q}
 \le \frac{1.002v}{64}\left(\frac23+\frac{1}{2Q}\right)
 < .71v/64.
\]
Here \(p\ge\sqrt{Z/q}\) and \(p\le\sqrt M\).
For the middle bound use the trapezoid inequality for the concave
square root and
\(\sqrt Z-\sqrt Q\ge\sqrt M(.9995-1/P)>\sqrt M/1.002\).
The lower bound on \(Z/M\) follows from
\(1-1/N-2H/N^2>.9995\), with \(H\le467\) below
\(20{,}000\) and \(H\le1131\) above.

The primes \(p\le P\) contribute at most
\(S(P)+\pi(P)V/N\). We have \(V<v+.001\) and
\[
 \frac{V}{\zeta}<3-\frac{3\log4-.001}{\zeta}.
\]
The displayed bound and its product with \(1+1.2762/\zeta\)
are increasing in \(\zeta\). At the upper corner of each row
they are bounded by \(A_v,A_c\), respectively.
In the last three rows both are less than three because
\(3\log4-.001>3\cdot1.2762\).
Combining the three ranges of primes, using \(S(P)<\log P-.49\)
and \(P/(P-1)<1.002\), gives the upper bound
\begin{equation}\label{eq:hard-upper}
 R=\zeta-.49+\frac PN\left(A_c+\frac{2(v+.001)}{64}\right)
 +\frac{.71v}{64}+\frac{2.004D\xi}{P}+\frac{2A_vD\xi}{N}.
\end{equation}

For the lower bound let \(E=D/b\).
The errors in \eqref{eq:pade-angles} satisfy
\(J<.883\sqrt\alpha\) and \(r/d<.641\sqrt\alpha\).
Indeed \(\arcsin x+x\sqrt{1-x^2}\le2x\) and
\(\arcsin x\le1.006x\) for \(x^2\le.03\).
For \(d\ge20\), use
\(\log C'>t-.6932-1.031E\),
\(g_0\le1.01+.128DE\), and \(g_0/d\le1.01/d+.128E\).
At fixed \(b\), the functions
\((1-1/d)(\log(b/d)-.6932)\) and
\([2\log(2d/\pi)-1.01]/d\) decrease as \(d\) increases
from the chosen \(d\ge20\) to \(D\). For the first,
use \(t<U<15\) at the chosen \(d\).
Thus \(\log\Lambda\) exceeds
\begin{equation}\label{eq:hard-lower}
 \begin{aligned}
 L={}&U-.001+(1-1/D)(t_D-.6932)
     +\frac{2\log(2D/\pi)-1.01}{D}\\
     &-1.159E-\bigl[.641(1.7032+.128DE)+.883\bigr]\sqrt E .
 \end{aligned}
\end{equation}
For \(d=2,3,\ 4\le d\le9,\ 10\le d\le19\), use
\(D=2,3,9,19\) respectively in \eqref{eq:hard-upper} and
\(t_D\), keeping the row's \(A_v,A_c\).
The same angular estimates used above give
\(L=U+ft_D-c\) with
\((f,c)=(1/2,.2),(2/3,.435),(.75,.68),(.9,.925)\),
respectively. For \(10\le d\le19\) retain the positive term
\(2\log(2d/\pi)/d\ge2\log(38/\pi)/19\) in
\eqref{eq:pade-angles}.

At fixed \(N\), \(L\) is concave in \(w\).
The bound \(R\) is convex in \(w\) provided
\[
 (P^2/N)\left(A_c+\frac{2(v+.001+6)}{64}\right)
      \ge 2.004D(6-\xi).
\]
This follows by differentiating the \(P/N\) and \(\xi/P\)
terms twice. The left side increases with \(U,w\), while
the right side decreases. It suffices to check \(w=0\) at the
lower \(N\) endpoint in each row for its largest \(D\).
At fixed \(w\), \(L\) grows with \(U\) at rate at least \(1.5\),
whereas \(R\) grows at rate at most \(2/3+2\cdot .71/64\).
The digit term with \(P/N\) and the two equal-height terms
decrease since \(\xi>3\).
The supplement checks the convexity inequality and the surplus
at the two endpoints of \(w\) at the smallest \(N\) in each row.
For \(d\ge20\) the certified margins are given in the table.
For all four smaller-gap cases the margin exceeds \(3.8\).
It follows that \(L>R\) throughout each row, which contradicts
\eqref{eq:pade-norm} and completes the proposition.

